
\section{Phase~IV Summary: Projective Cosmology, Structure, and the Topological Dark Sector}

Phase~IV extends the SAT framework into full cosmological geometry. 
We replace strain-based dynamics with projective geometry: cosmic expansion, inflation, and late-time acceleration emerge from the resolution of a filamentary manifold by an expanding time surface $\Sigma_t$. Structure arises from angular resistance and link intersections, while the dark sector emerges naturally from uncoupled or topologically inert bundles.

\tableofcontents

%-------------------------------------------------
\section{SAT Cosmology via Projection Geometry}

Assume a statistically isotropic ensemble of radially organized filaments $F(x)$.
The SAT cosmological model arises from a projective process:
a surface $\Sigma_t$ expands outward from the origin, intersecting structural filaments and projecting geometry into resolution.

The Hubble expansion rate follows directly from geometric resistance:
\[
H_{\text{SAT}}(t) = \frac{1}{R(t)} \cdot \frac{1}{\mathcal{R}(R(t))},
\]
where $\mathcal{R}(r)$ depends on:
\begin{itemize}
  \item $\rho_{\text{link}}(r)$ — filament intersection density
  \item $\theta_4(r)$ — local projection angle
  \item $\tau(r)$ — torsional suppression
\end{itemize}

This reproduces an FLRW-like metric in the large-scale average:
\[
ds^2 = -dt^2 + R(t)^2 d\vec{x}^2,
\]
with expansion arising not from vacuum energy, but from the evolving ease of resolution.

\subsection{Late-Time Acceleration and $\rho_\Lambda$}
Late-time acceleration is a natural consequence of diminishing intersection density and angular resistance:
\[
\lim_{t \to \infty} \mathcal{R}(R(t)) \to \text{const.} \quad \Rightarrow \quad H_{\text{SAT}}(t) \to \text{asymptote}.
\]
There is no need for a cosmological constant. $\rho_\Lambda$ is reinterpreted as a geometric saturation effect.

\subsection{Inflation from Projective Spike}
Initial high angular variance and dense filament intersection yield rapid early expansion. This is equivalent to inflation but emerges from projection mechanics, with no inflaton or tuning:
\[
W(R) = \int_0^R \rho_{\text{link}}(r) \cdot \theta_4(r) \cdot 4\pi r^2 dr \quad \text{peaks early.}
\]

%-------------------------------------------------
\section{Structure Formation and Angular Perturbations}

\subsection{Sources of Perturbation}
Fluctuations arise naturally from topological variations:
\begin{itemize}
  \item Local fluctuations in $\theta_4(x)$ and $\tau(x)$
  \item Angular misalignment between adjacent filaments
  \item Resolution interference across intersecting bundles
\end{itemize}

\subsection{Spectrum and Shape}
Projection fluctuations imprint on Σₜ as curvature modulations:
\[
P(k) \sim k^{n_s - 1}, \quad n_s \approx 0.96.
\]
Tensor modes emerge only from extreme τ-dominated bundles.

\subsection{Topological Non-Gaussianity}
Linked bundles and multi-filament entanglements generate:
\begin{itemize}
  \item Mild non-Gaussian signatures
  \item Possible anomalous features in the CMB
  \item Filament-bound topological defects at large scales
\end{itemize}

%-------------------------------------------------
\section{Dark Matter and Dark Energy Reinterpreted}

\subsection{Topological Sector Classification}
\begin{center}
\begin{tabular}{|c|c|c|c|c|}
\hline
Sector & Q & Phase State & Coupling & Behavior \\\hline
Visible & 1--3 & Locked & Electroweak/color & SM matter \\
Dark Matter & $\geq$1 & Incoherent & Gravity only & Cold, pressureless \\
Dark Energy & $\geq$3 & Frozen, high-$\tau$ & None & $\rho=\text{const}$ \\\hline
\end{tabular}
\end{center}

\subsection{Dark Matter as Residual Projection Resistance}
Bundles that resist resolution due to high τ or misalignment behave as CDM:
\begin{itemize}
  \item Gravitate normally, but remain topologically inert
  \item Do not project onto the SM sector
  \item Persist as cold, invisible geometry
\end{itemize}

\subsection{Dark Energy from Projection Exhaustion}
Regions with no remaining intersections and saturated $W(R)$ no longer resist Σₜ:
\[
\rho_\Lambda \sim \text{projective asymptote}.
\]

\subsection{Predictions}
\begin{itemize}
  \item No new particles required (WIMPs, axions unnecessary)
  \item Equation of state from geometry: $w \approx -1$
  \item Early acceleration and late flattening follow from angular projection profile
\end{itemize}

%-------------------------------------------------
\section{Conclusion}

Phase~IV shows that SAT, via PGCU, can reproduce cosmic expansion, structure emergence, and the full dark sector from projection geometry alone. This occurs without extra fields, vacuum energy, or fine-tuned potentials. The results unify time, mass, and curvature as emergent properties of Σₜ sweeping through a filamented topological manifold.
